\chapter{Continua i ich struktura}
\label{ch:topog25-continua}

\noindent\textbf{Cel rozdziału.}
Zwartość i spójność pozwalają wyróżnić przestrzenie, które
mogą mieć złożoną strukturę mimo braku rozdzielenia.
W rozdziale~\ref{ch:topog-spojnosc} omówiliśmy już
podstawowe własności spójności, składowe spójne,
spójność drogową i lokalną oraz krzywą sinusową topologa.
Nie powtarzamy tych dowodów.
Zbadamy istnienie minimalnych podcontinuów,
różnicę między continuami rozkładalnymi i nierozkładalnymi,
granice odwrotne i przykłady, w których lokalna spójność
zawodzi.
Skonstruujemy continuum nierozkładalne Knastera
oraz solenoid dyadyczny.
Na koniec wskażemy dokładny zakres związków
z funkcjami Peana.

\section{Continua i podcontinua}
\label{sec:topog25-basic}

\begin{definicja}[Continuum i podcontinuum]
\label{def:topog25-continuum}
\emph{Continuum} nazywamy niepustą, zwartą,
spójną przestrzeń Hausdorffa.
Jeżeli \(X\) jest continuum,
to \emph{podcontinuum} \(X\)
jest niepustym, spójnym, zwartym podzbiorem \(K\subseteq X\)
z topologią podprzestrzeni.
Continuum nazywamy \emph{nietrywialnym},
gdy zawiera co najmniej dwa punkty.
\emph{Continuum metryczne} to continuum,
którego topologia jest metryzowalna.
\end{definicja}

\begin{uwaga}[Rola aksjomatów]
\label{rem:topog25-closed}
W przestrzeni Hausdorffa każdy podzbiór zwarty
jest domknięty. Zatem podcontinua są zbiorami
domkniętymi w \(X\). Odwrotnie, każdy
niepusty, domknięty i spójny podzbiór
continuum \(X\) jest podcontinuum.
Definicja dopuszcza przestrzeń jednopunktową;
w niektórych zastosowaniach zastrzega się
nietrywialność, dlatego będziemy
ją jawnie zaznaczać, kiedy jest potrzebna.
\end{uwaga}

\begin{stwierdzenie}[Obrazy i produkty continuów]
\label{prop:topog25-images-products}
Niech \(X\) będzie continuum, a \(Y\)
przestrzenią Hausdorffa.
Jeżeli \(f\colon X\to Y\) jest ciągłe,
to \(f(X)\) jest continuum z topologią podprzestrzeni.
Ponadto iloczyn dowolnej niepustej rodziny
continuów jest continuum.
\end{stwierdzenie}
\begin{proof}
Obraz niepustego zbioru jest niepusty.
Jest zwarty jako ciągły obraz zwartego \(X\).
Jest spójny na mocy
Twierdzenia~\ref{thm:topog-spojnosc-obraz-ciagly};
jako podprzestrzeń \(Y\) jest Hausdorffa.
Stąd \(f(X)\) jest continuum.
Dla rodziny \((X_i)_{i\in I}\)
każdy czynnik jest niepusty i zwarty,
więc iloczyn jest niepusty i zwarty
na mocy twierdzenia Tychonoffa
(przy użyciu aksjomatu wyboru).
Jest Hausdorffa, a spójność produktu
wynika z
Twierdzenia~\ref{thm:topog-spojnosc-iloczyn}.
\end{proof}

\begin{przyklad}[Proste continua]
\label{ex:topog25-basic}
Nietrywialny odcinek \([a,b]\subseteq\mathbb R\)
jest continuum metrycznym.
Okrąg \(S^1\subseteq\mathbb R^2\)
jest continuum, ponieważ jest obrazem
odcinka przy odwzorowaniu
\(t\mapsto(\cos(2\pi t),\sin(2\pi t))\).
Z kolei zbiór Cantora jest zwarty,
ale nie jest continuum:
jest niespójny i zawiera więcej niż jeden punkt.
\end{przyklad}

\begin{twierdzenie}[Przecięcie zstępującego ciągu podcontinuów]
\label{thm:topog25-nested}
Niech \(X\) będzie zwartą przestrzenią Hausdorffa,
a \((K_n)_{n\ge1}\) ciągiem niepustych
zwartych i spójnych podzbiorów \(X\),
spełniającym \(K_{n+1}\subseteq K_n\).
Wówczas
\[
 K=\bigcap_{n=1}^{\infty}K_n
\]
jest niepustym podcontinuum \(X\).
\end{twierdzenie}
\begin{proof}
Zbiory \(K_n\) są domknięte w \(X\).
Mają własność skończonych przecięć,
bo każda skończona rodzina przecina się
w najmniejszym z rozważanych zbiorów.
Zwartość \(X\) daje \(K\neq\varnothing\).
Przecięcie zbiorów domkniętych
jest domknięte, zatem \(K\)
jest także zwarte.

Załóżmy, że \(K\) nie jest spójny.
Istnieją niepuste rozłączne zbiory
\(A,B\subseteq K\), otwarte i domknięte w \(K\),
takie że \(K=A\cup B\).
Są one zwarte, a więc domknięte w \(X\).
Normalność zwartego Hausdorffa \(X\)
daje rozłączne otwarte \(U,V\subseteq X\)
z \(A\subseteq U\), \(B\subseteq V\).
Wtedy \(K\subseteq U\cup V\).

Twierdzimy, że istnieje \(N\)
takie, że \(K_N\subseteq U\cup V\).
W przeciwnym razie niepuste
zwarte zbiory \(K_n\setminus(U\cup V)\)
tworzyłyby ciąg zstępujący.
Zwartość \(X\) dałaby punkt ich
przecięcia należący jednocześnie
do \(K\) i do \(X\setminus(U\cup V)\),
co przeczyłoby \(K\subseteq U\cup V\).
Zatem taki \(N\) istnieje.

Ponieważ \(A,B\subseteq K\subseteq K_N\),
oba zbiory \(K_N\cap U\) i \(K_N\cap V\)
są niepuste, otwarte w \(K_N\),
rozłączne i pokrywają \(K_N\).
Jest to rozdzielenie \(K_N\),
wbrew jego spójności.
\end{proof}

\begin{lemat}[Przecięcie skierowanej rodziny continuów]
\label{lem:topog25-chain-intersection}
Niech \(X\) będzie continuum
i \(\mathcal C\) niepustą rodziną jego
podcontinuów uporządkowaną liniowo
przez inkluzję.
Wówczas \(\bigcap_{C\in\mathcal C}C\)
jest niepustym podcontinuum \(X\).
\end{lemat}
\begin{proof}
Każdy \(C\in\mathcal C\) jest
niepusty i domknięty w zwartym \(X\).
Dzięki liniowemu uporządkowaniu
dowolne skończenie wiele członów
rodziny ma niepuste przecięcie,
równe najmniejszemu z nich.
Zatem \(K=\bigcap_{C\in\mathcal C}C\)
jest niepusty i zwarty.

Jeśli \(K\) byłby niespójny,
jak w dowodzie
Twierdzenia~\ref{thm:topog25-nested}
znaleźlibyśmy rozłączne otwarte
\(U,V\subseteq X\), których suma
zawiera \(K\) i które spotykają \(K\).
Zwartość \(X\) implikuje, że pewna
skończona podrodzina \(\mathcal C\)
ma przecięcie zawarte w \(U\cup V\):
istotnie, otwarte zbiory
\(X\setminus C\), \(C\in\mathcal C\),
wraz z \(U\cup V\) pokrywają \(X\);
wybieramy skończone podpokrycie.
Ponieważ \(\mathcal C\) jest łańcuchem,
jego przecięcie skończone jest
pewnym \(C_0\in\mathcal C\).
Wobec tego \(C_0\subseteq U\cup V\).
Ponadto \(K\subseteq C_0\)
spotyka zarówno \(U\), jak i \(V\).
Otrzymujemy rozdzielenie \(C_0\),
co jest sprzeczne ze spójnością.
\end{proof}

\section{Continua nierozkładalne i nieprzywiedlne}
\label{sec:topog25-indecomposable}

\begin{definicja}[Rozkładalność]
\label{def:topog25-indecomposable}
Continuum \(X\) jest
\emph{rozkładalne}, jeżeli istnieją
jego dwa właściwe podcontinua \(A,B\subsetneq X\)
takie, że \(X=A\cup B\).
W przeciwnym razie nazywamy \(X\)
\emph{nierozkładalnym}.
Nie należy utożsamiać rozkładalności
z niespójnością: continuum z definicji
jest spójne.
\end{definicja}

\begin{przyklad}[Przykłady rozkładalne]
\label{ex:topog25-decomposable}
Odcinek \([0,1]\) jest rozkładalny, gdyż
\([0,1]=[0,1/2]\cup[1/2,1]\),
a obie części są właściwymi podcontinuami.
Okrąg jest rozkładalny:
można przedstawić go jako sumę
dwóch domkniętych półokręgów.
Continuum jednoelementowe jest
nierozkładalne w przyjętej definicji,
ponieważ nie ma właściwego niepustego
podcontinuum; dlatego ważne przykłady
nierozkładalne będziemy wybierać
nietrywialne.
\end{przyklad}

\begin{stwierdzenie}[Właściwe podcontinua continuum nierozkładalnego]
\label{prop:topog25-nowhere-dense}
Niech \(X\) będzie nietrywialnym continuum,
w którym każde właściwe podcontinuum
ma puste wnętrze w \(X\).
Wówczas \(X\) jest nierozkładalne.
Równoważnie: jeżeli \(X\) jest rozkładalne,
to ma właściwe podcontinuum
o niepustym wnętrzu.
\end{stwierdzenie}
\begin{proof}
Gdyby \(X=A\cup B\) dla właściwych
podcontinuów \(A,B\), to \(B\)
byłby domknięty w \(X\), a zbiór
\(X\setminus B\) byłby niepusty i otwarty.
Z równości \(X=A\cup B\)
wynika \(X\setminus B\subseteq A\).
Zatem \(A\) byłoby właściwym
podcontinuum o niepustym wnętrzu,
wbrew założeniu.
\end{proof}

\begin{definicja}[Continuum nieprzywiedlne między punktami]
\label{def:topog25-irreducible}
Niech \(X\) będzie continuum,
a \(a,b\in X\).
Mówimy, że \(X\) jest
\emph{nieprzywiedlne między \(a\) i \(b\)},
jeśli nie istnieje właściwe
podcontinuum \(K\subsetneq X\)
zawierające jednocześnie \(a\) i \(b\).
\end{definicja}

\begin{twierdzenie}[Istnienie minimalnego continuum między punktami]
\label{thm:topog25-irreducible-existence}
Niech \(X\) będzie continuum,
a \(a,b\in X\).
Istnieje podcontinuum
\(K\subseteq X\), zawierające \(a,b\),
które jest nieprzywiedlne między
tymi punktami.
\end{twierdzenie}
\begin{proof}
Niech \(\mathcal P\) będzie rodziną
wszystkich podcontinuów \(C\subseteq X\)
zawierających \(a,b\).
Jest ona niepusta, ponieważ \(X\in\mathcal P\).
Uporządkujmy ją przez odwrotną inkluzję:
\(C\preceq D\) wtedy i tylko wtedy,
gdy \(D\subseteq C\).

Niech \(\mathcal L\subseteq\mathcal P\)
będzie niepustym łańcuchem.
Na mocy
Lematu~\ref{lem:topog25-chain-intersection}
zbiór
\(K_{\mathcal L}=\bigcap_{C\in\mathcal L}C\)
jest niepustym podcontinuum.
Zawiera \(a,b\) i jest górnym
ograniczeniem łańcucha
w porządku \(\preceq\).
Dla łańcucha pustego można
wybrać górne ograniczenie \(X\).
Lemat Kuratowskiego--Zorna
daje element maksymalny \(K\)
dla \(\preceq\), czyli
element minimalny dla inkluzji.
Gdyby właściwe podcontinuum
\(D\subsetneq K\) zawierało \(a,b\),
to \(D\in\mathcal P\) i \(K\prec D\),
wbrew maksymalności \(K\).
Zatem \(K\) jest nieprzywiedlne
między \(a,b\).
\end{proof}

\begin{uwaga}[Dwa różne pojęcia]
\label{rem:topog25-irred-vs-indec}
Nieprzywiedlność między punktami
i nierozkładalność nie są synonimami.
Odcinek \([0,1]\) jest
nieprzywiedlny między \(0\) i \(1\):
każdy spójny podzbiór prostej
zawierający te punkty zawiera
cały odcinek. Jednocześnie
odcinek jest rozkładalny.
\end{uwaga}

\section{Granice odwrotne przestrzeni zwartych}
\label{sec:topog25-inverse}

\begin{definicja}[Granica odwrotna ciągu]
\label{def:topog25-inverse-limit}
Niech \((X_n)_{n\ge0}\) będzie
ciągiem niepustych przestrzeni
zwartych Hausdorffa, a
\(f_n\colon X_{n+1}\to X_n\)
ciągłymi odwzorowaniami.
\emph{Granicę odwrotną} układu
\((X_n,f_n)\) definiujemy jako
\[
 \varprojlim(X_n,f_n)=
 \left\{(x_0,x_1,\ldots)\in
  \prod_{n=0}^{\infty}X_n:
  f_n(x_{n+1})=x_n\quad(n\ge0)\right\}
\]
z topologią podprzestrzeni
produktu. Ograniczenie \(n\)-tego
rzutowania do granicy oznaczamy
przez \(p_n\).
\end{definicja}

\begin{twierdzenie}[Granica odwrotna continuów]
\label{thm:topog25-inverse-connected}
Niech \((X_n)_{n\ge0}\) będzie
ciągiem continuów, a
\(f_n\colon X_{n+1}\to X_n\)
ciągłymi surjekcjami.
Wówczas granica odwrotna
\(X=\varprojlim(X_n,f_n)\)
jest continuum.
Wszystkie rzutowania
\(p_n\colon X\to X_n\)
są ciągłe i surjektywne.
Jeśli ponadto każde \(X_n\)
jest metryzowalne,
to \(X\) jest continuum metrycznym.
\end{twierdzenie}
\begin{proof}
Produkt \(P=\prod_{n\ge0}X_n\)
jest zwarty Hausdorffa na mocy
twierdzenia Tychonoffa.
Dla każdego \(n\) warunek
\(f_n(x_{n+1})=x_n\)
wyznacza zbiór domknięty w \(P\),
ponieważ jest przeciwobrazem
przekątnej domkniętej w
\(X_n\times X_n\).
Zatem \(X\) jest domknięty
w \(P\), a więc zwarty
i Hausdorffa.

Udowodnimy niepustość
oraz surjektywność \(p_n\).
Dla ustalonego \(n\) i
\(a\in X_n\) możemy wyznaczyć
kolejno \(x_{n+1},x_{n+2},\ldots\)
tak, aby
\(f_j(x_{j+1})=x_j\)
i \(x_n=a\); używamy kolejno
surjektywności \(f_j\)
oraz aksjomatu wyboru zależnego.
Dla \(j<n\) kładziemy
\(x_j=f_j\circ\cdots\circ
f_{n-1}(a)\).
Otrzymujemy punkt \(x\in X\)
z \(p_n(x)=a\).
W szczególności \(X\neq\varnothing\).

Pozostaje spójność.
Załóżmy przeciwnie, że
\(X=A\cup B\), gdzie \(A,B\)
są niepustymi rozłącznymi
zbiorami otwarto-domkniętymi
w \(X\). Oba są zwarte.
Pokażemy, że dla pewnego \(n\)
zwarte zbiory \(p_n(A)\)
i \(p_n(B)\) są rozłączne.
Gdyby dla każdego \(n\)
ich przecięcie było niepuste,
zbiór
\[
 F_n=\{(a,b)\in A\times B:p_n(a)=p_n(b)\}
\]
byłby niepusty i domknięty
w zwartym \(A\times B\).
Wynika z warunków zgodności,
że \(F_{n+1}\subseteq F_n\).
Zwartość dałaby parę
\((a,b)\in\bigcap_{n\ge0}F_n\).
Wtedy \(p_n(a)=p_n(b)\)
dla wszystkich \(n\), więc
\(a=b\) jako ciągi współrzędnych,
wbrew \(A\cap B=\varnothing\).
Dla pewnego \(n\)
mamy więc \(p_n(A)\cap p_n(B)=\varnothing\).

Ponieważ \(p_n\) jest surjektywne
oraz \(X=A\cup B\), otrzymujemy
\[
 X_n=p_n(A)\cup p_n(B).
\]
Oba zbiory po prawej są
niepuste, zwarte, a więc
domknięte w \(X_n\), i rozłączne.
Stanowią rozdzielenie \(X_n\),
sprzeczne ze spójnością \(X_n\).
To dowodzi spójności \(X\).

Jeśli każde \(X_n\) jest
metryzowalne, wybieramy zgodne
ograniczone metryki \(d_n\)
na \(X_n\). Wówczas
\[
 d(x,y)=\sum_{n=0}^{\infty}
  2^{-(n+1)}\min\{1,d_n(x_n,y_n)\}
\]
jest metryką na \(P\),
wyznaczającą topologię produktu.
Jej ograniczenie do \(X\)
metryzuje \(X\).
\end{proof}

\begin{stwierdzenie}[Continuum z rzutami przedziałowymi]
\label{prop:topog25-inverse-intervals}
Niech każde \(X_n=[0,1]\),
a \(f_n\colon[0,1]\to[0,1]\)
będzie ciągłą surjekcją.
Wówczas granica odwrotna
jest niepustym continuum metrycznym,
a wszystkie jej rzutowania
na przedział są surjektywne.
\end{stwierdzenie}
\begin{proof}
Przedział \([0,1]\)
jest nietrywialnym continuum
metrycznym; wystarczy zastosować
Twierdzenie~\ref{thm:topog25-inverse-connected}.
\end{proof}

\section{Continuum nierozkładalne Knastera}
\label{sec:topog25-knaster}

\begin{definicja}[Odwzorowanie namiotowe]
\label{def:topog25-tent}
Niech \(T\colon[0,1]\to[0,1]\)
będzie odwzorowaniem ciągłym danym wzorem
\[
 T(t)=
 \begin{cases}
  2t,&0\le t\le\tfrac12,\\
  2-2t,&\tfrac12\le t\le1.
 \end{cases}
\]
Przez \(K_T\) oznaczamy granicę
odwrotną stałego układu
\([0,1]\xleftarrow{T}[0,1]
\xleftarrow{T}\cdots\).
\end{definicja}

\begin{lemat}[Pełne gałęzie iteracji namiotu]
\label{lem:topog25-tent-branches}
Dla każdego \(r\ge1\) i
\(j\in\{0,\ldots,2^r-1\}\)
ograniczenie \(T^r\) do przedziału
\[
 I_{r,j}=\left[\frac{j}{2^r},
                     \frac{j+1}{2^r}\right]
\]
jest afiniczną bijekcją tego
przedziału na \([0,1]\).
W szczególności, dla każdego
niepustego otwartego podzbioru
\(U\subseteq[0,1]\), zbiór
\((T^r)^{-1}(U)\) spotyka
przedziały \(I_{r,0}\) i
\(I_{r,2^r-1}\).
\end{lemat}
\begin{proof}
Dla \(r=1\) stwierdzenie
wynika z definicji \(T\):
na obu połówkach przedziału
jest ono afiniczne i odwzorowuje
je bijektywnie na \([0,1]\).

Załóżmy tezę dla \(r\).
Na każdym \(I_{r,j}\)
odwzorowanie \(T^r\) jest
afiniczną bijekcją na \([0,1]\).
Przeciwobrazy dwóch przedziałów
\([0,\tfrac12]\) i
\([\tfrac12,1]\)
przy tym ograniczeniu dzielą
\(I_{r,j}\) na dwie równe połowy.
Na każdej z nich \(T^{r+1}=T\circ T^r\)
jest afiniczną bijekcją na \([0,1]\).
Są to przedziały
\(I_{r+1,2j}\) i \(I_{r+1,2j+1}\).
Indukcja dowodzi pierwszej tezy.
Druga wynika z surjektywności
każdego z tych ograniczeń.
\end{proof}

\begin{twierdzenie}[Continuum Knastera jest nierozkładalne]
\label{thm:topog25-knaster}
Continuum metryczne \(K_T\)
z Definicji~\ref{def:topog25-tent}
jest nietrywialne i nierozkładalne.
Co więcej, każde jego właściwe
podcontinuum ma puste wnętrze
w \(K_T\).
\end{twierdzenie}
\begin{proof}
Odwzorowanie \(T\) jest ciągłą
surjekcją \([0,1]\) na siebie.
Z Twierdzenia~\ref{thm:topog25-inverse-connected}
wynika, że \(K_T\) jest
niepustym continuum metrycznym,
a rzutowanie \(p_0\colon K_T\to[0,1]\)
jest surjektywne. Stąd \(K_T\)
ma co najmniej dwa punkty.

Przypuśćmy, że właściwe
podcontinuum \(C\subsetneq K_T\)
ma niepuste wnętrze.
Topologia granicy odwrotnej
ma bazę cylindrów
\(p_k^{-1}(U)\), gdzie
\(k\ge0\), a \(U\subseteq[0,1]\)
jest otwarty: dowolne skończone
warunki na wcześniejsze współrzędne
można przenieść do najwyższego
wśród ich indeksów.
Istnieją więc \(k\ge0\)
i niepusty otwarty \(U\subseteq[0,1]\)
takie, że
\[
 p_k^{-1}(U)\subseteq C.
 \tag{*}
\]

Niech \(j\ge k\) będzie dowolny.
Wybierzmy \(m\ge j+2\)
oraz \(m\ge k+2\)
i oznaczmy \(r=m-k\).
Surjektywność rzutowania \(p_m\)
i warunki granicy odwrotnej
dają
\[
 p_m(p_k^{-1}(U))=(T^r)^{-1}(U).
\]
W szczególności z \((*)\)
otrzymujemy
\((T^r)^{-1}(U)\subseteq p_m(C)\).
Zbiór \(p_m(C)\)
jest spójnym zwartym podzbiorem
przedziału \([0,1]\),
a zatem jest domkniętym przedziałem.
Lemat~\ref{lem:topog25-tent-branches}
pokazuje, że zawiera on
punkt z \([0,2^{-r}]\)
i punkt z \([1-2^{-r},1]\).
Ponieważ \(r\ge2\),
wynika stąd
\[
 [\tfrac14,\tfrac34]\subseteq p_m(C).
\]
Bezpośrednio z definicji \(T\)
obliczamy
\[
 T([\tfrac14,\tfrac34])=[\tfrac12,1],
 \qquad
 T^2([\tfrac14,\tfrac34])=[0,1].
\]
Ponieważ \(m-j\ge2\) oraz
\(T([0,1])=[0,1]\),
otrzymujemy
\[
 p_j(C)=T^{m-j}(p_m(C))=[0,1].
\]
Punkt \(j\ge k\) był dowolny.
Dla \(j<k\) również
\(p_j(C)=T^{k-j}(p_k(C))=[0,1]\).
Zatem wszystkie rzutowania
\(p_j(C)\) są pełnymi przedziałami.

To implikuje \(C=K_T\).
Istotnie, gdyby \(x\in K_T\setminus C\),
otwartość dopełnienia \(C\)
i bazowość cylindrów dały
\(j\) i otwarty \(W\subseteq[0,1]\)
takie, że
\(x\in p_j^{-1}(W)\subseteq K_T\setminus C\).
Wówczas \(p_j(C)\cap W=\varnothing\),
wbrew \(p_j(C)=[0,1]\).
Otrzymaliśmy sprzeczność.

Każde właściwe podcontinuum
\(K_T\) ma więc puste wnętrze.
Stwierdzenie~\ref{prop:topog25-nowhere-dense}
daje nierozkładalność \(K_T\).
\end{proof}

\begin{figure}[htbp]
\centering
\begin{tikzpicture}[x=5.0cm,y=3.0cm]
  \draw[->] (-.05,0)--(1.09,0) node[below] {$t$};
  \draw[->] (0,-.05)--(0,1.13) node[left] {$T(t)$};
  \draw[very thick,blue!65!black]
    (0,0)--(.5,1)--(1,0);
  \draw[densely dashed,gray] (.5,0)--(.5,1);
  \node[below,font=\small] at (0,0) {$0$};
  \node[below,font=\small] at (.5,0) {$\tfrac12$};
  \node[below,font=\small] at (1,0) {$1$};
  \node[left,font=\small] at (0,1) {$1$};
  \node[anchor=west,font=\small,align=left] at (.57,.88)
    {Każda z dwóch gałęzi\\pokrywa cały przedział};
\end{tikzpicture}
\caption{Odwzorowanie namiotowe \(T\); jego kolejne
iteracje mają \(2^r\) pełnych gałęzi afinicznych.
Ta własność jest kluczowa dla nierozkładalności \(K_T\).}
\label{fig:topog25-tent}
\end{figure}

\section{Solenoid dyadyczny a spójność lokalna}
\label{sec:topog25-solenoid}

\begin{definicja}[Solenoid dyadyczny]
\label{def:topog25-solenoid}
Niech \(S^1=\{z\in\mathbb C:|z|=1\}\).
\emph{Solenoidem dyadycznym} nazywamy
granicę odwrotną
\[
 \Sigma_2=\varprojlim
 \bigl(S^1\xleftarrow{z\mapsto z^2}S^1
 \xleftarrow{z\mapsto z^2}\cdots\bigr).
\]
Punkty \(\Sigma_2\) są ciągami
\((z_0,z_1,\ldots)\) liczb z \(S^1\)
spełniającymi \(z_{n+1}^2=z_n\).
\end{definicja}

\begin{stwierdzenie}[Solenoid jest continuum]
\label{prop:topog25-solenoid-continuum}
Solenoid \(\Sigma_2\) jest
nietrywialnym continuum metrycznym.
Działanie przez mnożenie współrzędnych
\[
 (z_n)_{n\ge0}\cdot(w_n)_{n\ge0}
 =(z_nw_n)_{n\ge0}
\]
czyni go ponadto zwartą
przemienną grupą topologiczną Hausdorffa.
\end{stwierdzenie}
\begin{proof}
Okrąg jest nietrywialnym
continuum metrycznym,
a odwzorowanie \(z\mapsto z^2\)
jest ciągłą surjekcją.
Twierdzenie~\ref{thm:topog25-inverse-connected}
daje pierwszą tezę wraz
z nietrywialnością, ponieważ
rzutowanie \(p_0\) jest surjektywne.

Jeżeli \(z_{n+1}^2=z_n\)
i \(w_{n+1}^2=w_n\),
to \((z_{n+1}w_{n+1})^2=z_nw_n\),
więc iloczyn współrzędny
ponownie należy do \(\Sigma_2\).
Elementem neutralnym jest
ciąg samych jedynek,
a odwrotnością ciągu \((z_n)\)
jest \((z_n^{-1})\).
Łączność i przemienność wynikają
z mnożenia na okręgu.
Mnożenie i odwracanie
są ciągłe w topologii produktu,
więc są ciągłe także
na podprzestrzeni \(\Sigma_2\).
\end{proof}

\begin{lemat}[Lokalny model solenoidu]
\label{lem:topog25-solenoid-local}
Niech \(J\subseteq\mathbb R\)
będzie niepustym przedziałem otwartym
o długości mniejszej niż \(1\),
a \(U=\{e^{2\pi it}:t\in J\}\subseteq S^1\).
Wówczas przeciwobraz
\(p_0^{-1}(U)\subseteq\Sigma_2\)
jest homeomorficzny z
\(J\times\{0,1\}^{\mathbb N}\),
gdzie drugi czynnik ma topologię produktową.
\end{lemat}
\begin{proof}
Odwzorowanie
\(t\mapsto e^{2\pi it}\)
jest homeomorfizmem \(J\to U\),
ponieważ \(|J|<1\).
Dla \(t\in J\) oraz
\(\varepsilon=(\varepsilon_1,\varepsilon_2,\ldots)
\in\{0,1\}^{\mathbb N}\)
konstruujemy rekurencyjnie
\(t_0=t\) oraz
\[
 t_{n+1}=\frac{t_n+\varepsilon_{n+1}}{2},
 \qquad z_n=e^{2\pi it_n}.
\]
Wtedy \(z_{n+1}^2=z_n\),
więc \((z_n)\in p_0^{-1}(U)\).
Daje to odwzorowanie
\(\Phi\colon J\times\{0,1\}^{\mathbb N}
\to p_0^{-1}(U)\).
Każda współrzędna \(z_n\)
zależy ciągle od \(t\)
i od skończonej liczby
\(\varepsilon_1,\ldots,\varepsilon_n\),
zatem \(\Phi\) jest ciągłe.

Z danego ciągu \((z_n)\)
odtwarzamy jednoznacznie \(t_0\in J\).
W każdym kroku \(z_{n+1}\)
jest jednym z dwóch różnych
pierwiastków kwadratowych z \(z_n\),
którym odpowiadają dokładnie
dwie wartości \(t_{n+1}\)
w podanym wzorze.
Wybór \(\varepsilon_{n+1}\)
jest więc jednoznaczny
i otrzymujemy odwrotność \(\Phi\).

Aby wykazać ciągłość tej odwrotności,
ustalmy punkt i dowolne \(N\).
Po ograniczeniu \(t_0\)
do małego przedziału wokół
wybranej wartości oba pierwiastki
kwadratowe \(z_0\) można rozróżnić
przez rozłączne otwarte łuki.
To zapewnia lokalną stałość
\(\varepsilon_1\).
Następnie postępujemy analogicznie
dla kolejnych pierwiastków,
uzyskując ciągłość każdej
skończonej listy
\((\varepsilon_1,\ldots,\varepsilon_N)\).
Są to wszystkie podstawowe warunki
topologii produktowej
na \(\{0,1\}^{\mathbb N}\),
więc odwrotność \(\Phi^{-1}\)
jest ciągła.
\end{proof}

\begin{twierdzenie}[Solenoid nie jest lokalnie spójny]
\label{thm:topog25-solenoid-not-locally-connected}
Solenoid dyadyczny \(\Sigma_2\)
jest continuum metrycznym,
które nie jest lokalnie spójne
w żadnym punkcie.
\end{twierdzenie}
\begin{proof}
Niech \(x\in\Sigma_2\).
Wybierzmy krótki otwarty łuk
\(U\subsetneq S^1\) zawierający
pierwszą współrzędną \(p_0(x)\).
Z Lematu~\ref{lem:topog25-solenoid-local}
otwarte otoczenie
\(p_0^{-1}(U)\) punktu \(x\)
jest homeomorficzne
z \(J\times\{0,1\}^{\mathbb N}\).
Przestrzeń Cantora
\(\{0,1\}^{\mathbb N}\)
jest całkowicie niespójna
i nie ma punktów izolowanych.
Dlatego każdy spójny podzbiór
\(J\times\{0,1\}^{\mathbb N}\)
ma stałą drugą współrzędną:
jej obraz przy rzutowaniu
na przestrzeń Cantora jest spójny,
więc jest singletonem.

Przypuśćmy, że \(\Sigma_2\)
jest lokalnie spójny w \(x\).
Wówczas wewnątrz
\(p_0^{-1}(U)\) istnieje
spójne otwarte otoczenie \(W\ni x\).
Jego obraz przez \(\Phi^{-1}\)
musiałby leżeć w
\(J\times\{\varepsilon\}\)
dla pewnego ciągu \(\varepsilon\).
Jednocześnie byłby
niepusty i otwarty
w \(J\times\{0,1\}^{\mathbb N}\).
To niemożliwe, ponieważ
singleton \(\{\varepsilon\}\)
nie jest otwarty w przestrzeni Cantora.
Otrzymujemy sprzeczność.
Punkt \(x\) był dowolny.
\end{proof}

\section{Continua Peana i odwzorowania przedziału}
\label{sec:topog25-peano}

\begin{definicja}[Continuum Peana]
\label{def:topog25-peano}
\emph{Continuum Peana} nazywamy
continuum metryczne,
które jest lokalnie spójne.
Takie przestrzenie obejmują
odcinki, okręgi i skończone
spójne grafy topologiczne.
\end{definicja}

\begin{definicja}[Skończony graf topologiczny]
\label{def:topog25-finite-graph}
\emph{Skończonym grafem topologicznym}
nazywamy realizację geometryczną
skończonego grafu kombinatorycznego,
otrzymaną przez sklejenie
skończenie wielu egzemplarzy \([0,1]\)
w wierzchołkach, przy założeniu,
że wnętrza różnych krawędzi
są rozłączne.
Dopuszczamy graf jednoelementowy
bez krawędzi.
\end{definicja}

\begin{twierdzenie}[Skończone grafy spójne są obrazami przedziału]
\label{thm:topog25-finite-graph-peano}
Niech \(P\) będzie realizacją
geometryczną skończonego,
spójnego grafu kombinatorycznego.
Wówczas \(P\) jest continuum Peana.
Ponadto istnieje ciągła surjekcja
\(\gamma\colon[0,1]\to P\).
\end{twierdzenie}
\begin{proof}
Jeśli graf ma jeden wierzchołek
i nie ma krawędzi, teza jest oczywista:
wybieramy odwzorowanie stałe.
Dalej załóżmy, że graf ma
przynajmniej jedną krawędź.

Realizacja \(P\) jest ilorazem
skończonej sumy rozłącznej
odcinków, zidentyfikowanych
w skończenie wielu końcach.
Jest zwarta, ponieważ
jest ciągłym obrazem
skończonej sumy zwartych odcinków.
Można ją metryzować
metryką długości najkrótszej
drogi wzdłuż krawędzi
(po przypisaniu każdej krawędzi
długości \(1\)).
Dla różnych punktów odległość
jest dodatnia; trójkątna nierówność
wynika ze składania dróg.
Topologia tej metryki jest
topologią realizacji:
we wnętrzu krawędzi ma bazę
małych przedziałów, a w wierzchołku
ma bazę małych gwiazd
złożonych z początkowych odcinków
wszystkich incydentnych krawędzi.
Są to zbiory otwarte i spójne.
Zatem \(P\) jest lokalnie spójna.

Spójność grafu kombinatorycznego
oznacza, że dowolne dwa wierzchołki
łączy skończony ciąg krawędzi.
Dołączając początkowy i końcowy
odcinek krawędzi, otrzymujemy
drogę ciągłą między dowolnymi
punktami \(P\).
Stąd \(P\) jest spójna,
a więc jest continuum Peana.

Aby skonstruować surjekcję
z przedziału, wybierzmy wierzchołek
bazowy \(v\).
Ponieważ graf jest skończony i spójny,
dla każdej krawędzi \(e\)
istnieje skończona droga z \(v\)
do jednego końca \(e\).
Przechodzimy tę drogę,
następnie pokonujemy krawędź \(e\),
po czym wracamy do \(v\)
wzdłuż odwróconej drogi
z drugiego końca \(e\).
Otrzymujemy zamknięty spacer
zaczynający się i kończący w \(v\),
którego obraz zawiera całą krawędź \(e\).
Wykonujemy kolejno takie spacery
dla wszystkich krawędzi grafu.
Jest ich skończenie wiele,
a wszystkie łączą się w \(v\).
Parametryzując kolejne odcinki
spacerów liniowo na odpowiednich
podprzedziałach \([0,1]\),
otrzymujemy ciągłe odwzorowanie
\(\gamma\colon[0,1]\to P\).
Jego obraz zawiera każdą krawędź
oraz wszystkie wierzchołki,
więc jest równy \(P\).
\end{proof}

\begin{twierdzenie}[Istnienie krzywych wypełniających sześcian]
\label{thm:topog25-space-filling}
Dla każdej liczby całkowitej \(n\ge1\)
istnieje ciągła surjekcja
\[
 F_n\colon[0,1]\longrightarrow[0,1]^n.
\]
W szczególności dla \(n=2\)
istnieje krzywa Peana wypełniająca
cały kwadrat.
\end{twierdzenie}
\begin{proof}
Niech \(C\subseteq[0,1]\)
będzie klasycznym zbiorem Cantora.
Przypomnijmy z
Przykładu~\ref{ex:topog23-cantor},
że odwzorowanie
\[
 c\colon\{0,1\}^{\mathbb N}\longrightarrow C,
 \qquad c((\varepsilon_k)_{k\ge1})
 =\sum_{k=1}^{\infty}\frac{2\varepsilon_k}{3^k}
\]
jest homeomorfizmem.
Dla ustalonego \(n\ge1\)
definiujemy na przestrzeni
\(\{0,1\}^{\mathbb N}\)
funkcje
\[
 b_j((\varepsilon_k)_{k\ge1})
 =\sum_{m=1}^{\infty}
 \frac{\varepsilon_{n(m-1)+j}}{2^m},
 \qquad 1\le j\le n.
\]
Każda \(b_j\) jest ciągła:
suma pierwszych \(M\) wyrazów
zależy od skończenie wielu
współrzędnych, a ogon ma moduł
co najwyżej \(2^{-M}\),
jednostajnie względem wszystkich
ciągów \((\varepsilon_k)\).
Wobec tego
\(b=(b_1,\ldots,b_n)\)
jest ciągłym odwzorowaniem
do sześcianu \([0,1]^n\).

Odwzorowanie \(b\) jest surjektywne.
Istotnie, dla dowolnego
\(y=(y_1,\ldots,y_n)\in[0,1]^n\)
wybieramy rozwinięcia dwójkowe
\[
 y_j=\sum_{m=1}^{\infty}
       \frac{\eta_{j,m}}{2^m},
 \qquad \eta_{j,m}\in\{0,1\}.
\]
Dla liczby \(1\) wybieramy
rozwinięcie z samych jedynek.
Przypisując
\(\varepsilon_{n(m-1)+j}=\eta_{j,m}\),
otrzymujemy ciąg, dla którego
\(b((\varepsilon_k))=y\).
Zatem odwzorowanie
\(f=b\circ c^{-1}\colon C\to[0,1]^n\)
jest ciągłe i surjektywne.

Zbiór \(C\) jest domknięty
w normalnej przestrzeni \([0,1]\).
Każda współrzędna \(f_j\colon C\to[0,1]\)
ma więc, na mocy
twierdzenia Tietzego o przedłużeniu
\ref{thm:topog-tietze},
ciągłe przedłużenie
\(F_j\colon[0,1]\to[0,1]\).
Odwzorowanie
\(F_n=(F_1,\ldots,F_n)\colon[0,1]\to[0,1]^n\)
jest ciągłe.
Jego obraz zawiera \(F_n(C)=f(C)=[0,1]^n\),
a jednocześnie jest podzbiorem \([0,1]^n\);
wobec tego jest surjekcją.
\end{proof}

\begin{uwaga}[Zakres twierdzenia Hahna--Mazurkiewicza]
\label{rem:topog25-hahn-mazurkiewicz}
Klasyczne twierdzenie Hahna--Mazurkiewicza
mówi, że niepusta przestrzeń Hausdorffa
jest ciągłym obrazem przedziału \([0,1]\)
wtedy i tylko wtedy, gdy jest
zwarta, spójna, lokalnie spójna
i metryzowalna.
Pełny dowód obu implikacji
wymaga osobnej konstrukcji
kolejnych coraz drobniejszych
parametryzacji; nie wynika
z udowodnionego wyżej
twierdzenia o skończonych grafach.
W niniejszym rozdziale
podaliśmy pełny dowód
dla skończonych grafów topologicznych,
natomiast ogólną równoważność
zaznaczamy jako klasyczny
wynik do szczegółowego opracowania
w dalszej iteracji.
Solenoid z
Twierdzenia~\ref{thm:topog25-solenoid-not-locally-connected}
nie spełnia warunku lokalnej spójności,
mimo że jest continuum metrycznym.
\end{uwaga}

\par\smallskip\noindent\textbf{Podsumowanie.}
Minimalne podcontinuum między dowolnymi
dwoma punktami istnieje dzięki
lematowi Kuratowskiego--Zorna
i twierdzeniu o przecięciu łańcucha.
Granice odwrotne ciągów continuów
z surjektywnymi odwzorowaniami wiążącymi
pozostają continuami.
Dla odwzorowania namiotowego
otrzymaliśmy nietrywialne
continuum nierozkładalne Knastera.
Solenoid dyadyczny jest natomiast
continuum metrycznym i zwartą
grupą topologiczną, lecz
nie jest lokalnie spójny.
Dla skończonych grafów spójnych
i wszystkich sześcianów \\( [0,1]^n \\)
skonstruowaliśmy ciągłe surjekcje
z odcinka, bez powtarzania
podstaw spójności z rozdziału~08.